\SourceLeaf{61}
$\gamma$) Once these are known, finite expressions for the remaining coefficients are readily derived by the recurrence scheme found earlier. It was
\[
A^{(h+2)} = \frac{n(\nu-h)A^{(h)} - 2(h+1)A^{(h+1)}}{n(\nu+h+2)}
\]
and hence, in this case, because $\nu=-1$,
\[
A^{(h+2)} = \frac{-n(h+1)A^{(h)} - 2(h+1)A^{(h+1)}}{n(h+1)}
\]
or
\[
A^{(h+2)} = -\Bigl(A^{(h)} + \frac{2}{n}A^{(h+1)}\Bigr).
\]
If you carry out the calculations using this recurrence scheme, you will encounter a very simple law\HistoricalNote{*)}{*)}{The expression of this law in \textsc{Euler}, loc. cit., \S\,275 is incorrect, perhaps through an error by the typesetter.}, namely
\[
A^{(h)} = (-1)^h\,\frac{2}{\BareRoot{(1-nn)}}\left(\frac{1-\BareRoot{(1-nn)}}{n}\right)^{h}.
\]
To establish it by proof, suppose it holds for the coefficients up through the $h^{\text{tum}}$. By the recurrence scheme we shall have
\[
\begin{aligned}
A^{(h+1)} & = (-1)^h\,\frac{2}{\BareRoot{(1-nn)}}\left(\frac{1-\BareRoot{(1-nn)}}{n}\right)^{h-1}\\&\quad + (-1)^{h+1}\,\frac{2}{n}\,\frac{2}{\BareRoot{(1-nn)}}\left(\frac{1-\BareRoot{(1-nn)}}{n}\right)^{h}
\end{aligned}
\]
or, taking out the common factor,
\[
(-1)^h\,\frac{2}{\BareRoot{(1-nn)}}\left(\frac{1-\BareRoot{(1-nn)}}{n}\right)^{h-1}\left\{1-\frac{2}{n}\left(\frac{1-\BareRoot{(1-nn)}}{n}\right)\right\}.
\]
The factor enclosed in braces is
\[
\begin{aligned}
&= \frac{nn-2+2\BareRoot{(1-nn)}}{nn}\\&= -\frac{(1-nn)-2\BareRoot{(1-nn)}+1}{nn}\\&= -\left(\frac{1-\BareRoot{(1-nn)}}{n}\right)^{2},
\end{aligned}
\]
and, on including it, we obtain
\[
\begin{aligned}
A^{(h+1)} & = (-1)^h\,\frac{2}{\BareRoot{(1-nn)}}\left(\frac{1-\BareRoot{(1-nn)}}{n}\right)^{h-1}\\&\quad .  -\left(\frac{1-\BareRoot{(1-nn)}}{n}\right)^{2}\\&= (-1)^{h+1} . \frac{2}{\BareRoot{(1-nn)}}\left(\frac{1-\BareRoot{(1-nn)}}{n}\right)^{h+1},
\end{aligned}
\]
from which the general validity of the law is evident.

